% =============================================================================
\section{Sequential Models: LSTM and BiLSTM}
\label{sec:lstm}
% =============================================================================

\subsection{Recurrent Neural Networks}

To incorporate context into word representations, researchers turned to
\textbf{Recurrent Neural Networks (RNNs)}, which process a sequence of tokens
$x_1, x_2, \ldots, x_T$ one step at a time, maintaining a hidden state $h_t$
that accumulates information from previous tokens:
\begin{equation}
  h_t = f_\theta(h_{t-1},\, x_t)
\end{equation}
where $f_\theta$ is a learned function. In principle, $h_T$ encodes the entire
history of the sequence. In practice, simple RNNs suffer from
\emph{vanishing gradients}~\cite{bengio1994learning}: gradients shrink
exponentially as they are back-propagated through many timesteps, making it
impossible to learn long-range dependencies.

\subsection{Long Short-Term Memory (LSTM)}
\label{subsec:lstm}

The \textbf{Long Short-Term Memory} (LSTM)~\cite{hochreiter1997lstm}
(Hochreiter \& Schmidhuber, 1997) addresses the vanishing gradient problem by
introducing a \emph{cell state} $c_t$ alongside the hidden state, regulated by
three learned gates: the \emph{forget gate} $f_t$, the \emph{input gate} $i_t$,
and the \emph{output gate} $o_t$. The full update equations are:

\begin{align}
  f_t          &= \sigma(W_f [h_{t-1}, x_t] + b_f)             \label{eq:forget}    \\
  i_t          &= \sigma(W_i [h_{t-1}, x_t] + b_i)             \label{eq:inputgate} \\
  \tilde{c}_t  &= \tanh(W_c [h_{t-1}, x_t] + b_c)              \label{eq:candidate} \\
  c_t          &= f_t \odot c_{t-1} + i_t \odot \tilde{c}_t    \label{eq:cell}      \\
  o_t          &= \sigma(W_o [h_{t-1}, x_t] + b_o)             \label{eq:outgate}   \\
  h_t          &= o_t \odot \tanh(c_t)                          \label{eq:hidden}
\end{align}

\noindent
where $\sigma$ is the sigmoid function, $\odot$ denotes element-wise
multiplication, and $[h_{t-1}, x_t]$ denotes concatenation.

The cell state $c_t$ (Eq.~\ref{eq:cell}) acts as a \emph{memory conveyor belt}:
the forget gate decides what to discard from the previous state, and the input
gate controls what new information to write. This mechanism allows gradients to
flow across hundreds of timesteps, resolving the vanishing gradient problem that
plagued simple RNNs.

% ── LSTM encoder–decoder diagram ──────────────────────────────────────────────
\begin{figure}[H]
\centering
\begin{tikzpicture}[
  >=Stealth, thick,
  enc/.style={draw=green!65!black, fill=green!10, rounded corners=4pt,
              minimum width=0.85cm, minimum height=2.2cm, line width=1.2pt},
  dec/.style={draw=blue!75!black, fill=cyan!10, rounded corners=4pt,
              minimum width=0.85cm, minimum height=2.2cm, line width=1.5pt},
  emb/.style={draw=green!65!black, fill=yellow!12, rounded corners=3pt,
              minimum width=2.0cm, minimum height=0.55cm},
  rdot/.style={circle, fill=red!80, minimum size=5pt, inner sep=0pt},
  bdot/.style={circle, fill=blue!75, minimum size=5pt, inner sep=0pt},
  rarr/.style={-{Stealth[length=5pt,width=4pt]}, thick, red!80},
  barr/.style={-{Stealth[length=5pt,width=4pt]}, thick, blue!75},
  scale=0.82, transform shape
]

%% Language labels
\node[font=\small\bfseries, text=green!60!black] at (2.8,2.8)  {English (Encoder)};
\node[font=\small\bfseries, text=blue!70!black]  at (8.2,2.8)  {French (Decoder)};

%% Encoder cells
\foreach \x/\lab/\nm in {0.0/$h_1$/h1, 2.6/$h_2$/h2, 5.2/$h_3$/h3}{
  \node[enc](\nm) at (\x,0.2){};
  \foreach \y in {0.9, 0.45, 0.0, -0.45}{\node[rdot] at (\x,\y){};}
  \node[font=\small\itshape] at (\x,1.55) {\lab};
}
\draw[rarr] (-1.1,0.2) -- (h1.west);
\draw[rarr] (h1.east) -- (h2.west) node[midway,above,font=\scriptsize\itshape]{W};
\draw[rarr] (h2.east) -- (h3.west) node[midway,above,font=\scriptsize\itshape]{W};

%% Input embeddings
\foreach \xc/\lbl/\word in {0.0/$x_1$/I, 2.6/$x_2$/am, 5.2/$x_3$/happy}{
  \node[emb](xe\word) at (\xc,-1.7){};
  \foreach \dx in {-0.6,-0.3,0,0.3,0.6}{\node[rdot] at (\xc+\dx,-1.7){};}
  \node[font=\scriptsize\itshape] at (\xc-0.8,-1.25) {\lbl};
  \node[font=\small] at (\xc,-2.35) {\word};
  \draw[rarr] (xe\word.north) -- ++(0,0.45);
}
\draw[rarr] (xeI.north)++(0,0.45) -- (h1.south);
\draw[rarr] (xeam.north)++(0,0.45) -- (h2.south);
\draw[rarr] (xehappy.north)++(0,0.45) -- (h3.south);

%% Decoder cells
\node[dec](d1) at (7.2,0.2){};
\foreach \y in {0.9,0.45,0.0,-0.45}{\node[bdot] at (7.2,\y){};}
\node[dec](d2) at (9.0,0.2){};
\foreach \y in {0.9,0.45,0.0,-0.45}{\node[bdot] at (9.0,\y){};}
\node[font=\small\itshape, text=blue!70!black] at (7.2,-1.1) {$y_1$};
\node[font=\small\itshape, text=blue!70!black] at (9.0,-1.1) {$y_2$};
\node[font=\small\bfseries, text=blue!75!black] at (7.2,2.1) {Je};
\node[font=\small\bfseries, text=blue!75!black] at (9.0,2.1) {suis};
\draw[barr] (h3.east) -- ++(0.55,0) |- (d1.west);
\draw[barr] (d1.east) -- (d2.west);
\draw[barr] (d1.north) -- ++(0,0.5);
\draw[barr] (d2.north) -- ++(0,0.5);

%% Timestep annotations
\foreach \x/\t in {0.0/1, 2.6/2, 5.2/3}{
  \node[font=\normalsize\bfseries, text=teal!80!black] at (\x,-3.0) {$t=\t$};
}
\end{tikzpicture}
\caption{LSTM encoder--decoder for sequence-to-sequence translation
  (\textit{``I am happy''} $\to$ \textit{``Je suis \ldots''}). The encoder
  processes tokens sequentially, passing hidden state $h_t$ forward; the
  decoder generates the target sequence conditioned on the final encoder state.}
\label{fig:lstm}
\end{figure}

\subsection{Limitations of the LSTM}

While the LSTM addresses the vanishing gradient problem, it has three remaining
weaknesses that limit its effectiveness for language understanding:

\begin{description}[leftmargin=2em, style=nextline]
  \item[\textbf{1.~Sequential bottleneck.}]
    The recurrence $h_t = f(h_{t-1}, x_t)$ is inherently sequential: token $t$
    cannot be processed until token $t{-}1$ is complete. This prevents
    parallelisation along the time axis, making training slow on long sequences.

  \item[\textbf{2.~Memory degradation.}]
    Despite the gating mechanism, information from early tokens is progressively
    diluted over many timesteps. For a sequence of length $T$, the influence of
    $x_1$ on $h_T$ decays roughly as $\prod_{t=2}^{T} f_t$ --- a product of
    forget gate values all $\leq 1$.

  \item[\textbf{3.~Unidirectionality.}]
    A left-to-right LSTM builds the representation of token $t$ from only its
    left context $x_1, \ldots, x_{t-1}$. The right context --- which is often
    crucial for disambiguation --- is unavailable.
\end{description}

\subsection{Bidirectional LSTM (BiLSTM)}
\label{subsec:bilstm}

The \textbf{Bidirectional LSTM}~\cite{schuster1997bilstm} (BiLSTM) attempts to
address the unidirectionality problem by running two separate LSTM passes over
the input --- a forward pass (left-to-right) producing $\overrightarrow{h_t}$
and a backward pass (right-to-left) producing $\overleftarrow{h_t}$ --- and
concatenating their outputs:
\begin{equation}
  h_t = \left[\overrightarrow{h_t}\,;\,\overleftarrow{h_t}\right]
  \label{eq:bilstm}
\end{equation}

% ── BiLSTM diagram ────────────────────────────────────────────────────────────
\begin{figure}[H]
\centering
\begin{tikzpicture}[
  >=Stealth, x=1.6cm, y=1.2cm,
  xnode/.style={circle, draw=blue!70!black, fill=cyan!18,
                minimum size=8mm, inner sep=0pt, font=\small},
  ynode/.style={circle, draw=violet!75!black, fill=magenta!16,
                minimum size=8mm, inner sep=0pt, font=\small},
  fcell/.style={draw=teal!65!black, fill=teal!30, rounded corners=2pt,
                minimum width=9mm, minimum height=8mm, font=\small\bfseries},
  bcell/.style={draw=cyan!70!black, fill=cyan!22, rounded corners=2pt,
                minimum width=9mm, minimum height=8mm, font=\small\bfseries},
  fwd/.style={->, very thick, teal!85!black},
  bwd/.style={->, very thick, cyan!75!black},
  conn/.style={->, thick, black!70},
  scale=0.9, transform shape
]

\def\ymid{0} \def\ytop{1.1} \def\ybot{-1.1}

\foreach \x/\idx in {0/0,1/1,2/2,3/i}{
  \node[xnode](x\idx) at (\x,\ybot) {$x_{\idx}$};
  \node[ynode](y\idx) at (\x,\ytop+1.0) {$y_{\idx}$};
  \node[fcell](f\idx) at (\x,\ymid) {$A$};
  \node[bcell](b\idx) at (\x,\ytop) {$A'$};
  \draw[conn] (x\idx)--(f\idx);
  \draw[conn] (f\idx.north) to[out=90,in=-90] (b\idx.south);
  \draw[conn] (b\idx.north) -- (y\idx.south);
}
\node[font=\normalsize] at (2.55,\ybot) {$\cdots$};

\draw[fwd] (-0.6,\ymid) -- (f0);
\draw[fwd] (f0)--(f1) (f1)--(f2) (f2)--++(0.55,0);

\draw[bwd] (3.6,\ytop) -- (bi);
\draw[bwd] (bi)--(b2) (b2)--(b1) (b1)--(b0) (b0)--++(-0.55,0);

\node[draw=black!20, fill=white, rounded corners=2pt,
      font=\scriptsize, align=left, inner sep=4pt]
  at (1.5,-1.85)
  {$A$\,: forward cell \quad $A'$\,: backward cell \quad
   $x_t$\,: input \quad $y_t$\,: output};

\node[font=\small\bfseries] at (1.5, 2.8)
  {$h_t = [\overrightarrow{h_t}\,;\,\overleftarrow{h_t}]$ \ (concatenation)};

\end{tikzpicture}
\caption{Bidirectional LSTM (BiLSTM). A forward pass (teal, $A$) processes
  tokens left-to-right; a backward pass (cyan, $A'$) processes tokens
  right-to-left. The outputs are concatenated at each position to form a
  context-aware representation. The two passes are trained independently and
  merged only at the output.}
\label{fig:bilstm}
\end{figure}

The BiLSTM provides richer representations than a unidirectional LSTM.
ELMo~\cite{peters2018elmo} (Peters et al., 2018), a prominent pre-\bert{}
contextualised embedding model, uses a stacked BiLSTM and demonstrated that
contextualised representations transfer effectively to downstream tasks.
However, the BiLSTM still inherits two critical limitations:

\begin{itemize}[leftmargin=2em]
  \item \textbf{Not truly bidirectional.} The forward and backward passes are
        trained \emph{separately} and only combined at the output
        (Eq.~\ref{eq:bilstm}). The two directions never interact during feature
        extraction; the model cannot, within a single layer, allow left-context
        and right-context information to jointly influence each other.
  \item \textbf{Still sequential.} Both passes remain sequential along their
        respective directions, limiting parallelisation and scalability.
\end{itemize}
